% --- Archon physics formalization source begin ---
% archon:physics
% archon:covers IPhO2026Problems/problem_IPhO_2026_2_C_1.lean
% archon:source-report reports/ipho_2026/problem_IPhO_2026_2_C_1.source.json
% archon:problem-id IPhO_2026_2
% archon:part-id C.1

\chapter{Physics problem IPhO\_2026\_2\_C\_1}
\label{ch:IPhO2026Problems_problem_IPhO_2026_2_C_1}

\paragraph{Problem source.}
For the half-cylindrical mirror of radius R, ray A is incident at angle theta
and its reflected line is y = m\_A*x + b\_A.  A neighboring parallel ray B is
incident at theta + Delta theta, with Delta theta much smaller than theta, and
its reflected line is y = m\_B*x + b\_B.  The envelope/intersection of neighboring
rays forms the caustic.  Use Figure 2g and its coordinate convention.

Current subquestion:
Write the slope m\_A and intercept b\_A of reflected ray A in terms of theta and R.

\paragraph{Current subquestion.}
Write the slope m\_A and intercept b\_A of reflected ray A in terms of theta and R.

\paragraph{Recorded answer/context.}
m\_A = cot(2*theta), and b\_A = R/(2*cos(theta)).

\paragraph{Figure/image path.}
/root/proposal\_for\_physic/science-mango/ipho\_2026\_source/image/T2\_page-4.png

\paragraph{Formalization target.}
create a compiling Lean file with sorry bodies at `IPhO2026Problems/problem\_IPhO\_2026\_2\_C\_1.lean`. The Lean declarations must preserve the physical quantities, dimensions or dimensional roles, figure labels, governing-law hypotheses, and final relation expressed by this problem.
Use Mathlib/Physlib names found through LeanExplore where available. If a physics API is missing, introduce faithful local abstractions rather than scalar placeholder aliases.

\paragraph{Iteration 002 redraft contract.}
Import Physlib and make the mirror radius, point coordinates, and line
intercept length-valued.  If the analytic geometry is written in scalar
coordinates, expose a named SI-length projection and state every coordinate
and intercept equation through it.  The slope, angles, and normalized
direction components are dimensionless.  Preserve the upward incoming ray,
radial outward normal, down-left reflected branch, and the vector
specular-reflection equation.

\begin{theorem}[Physics formalization target]
  \label{thm:physics:IPhO_2026_2_C_1:target}
  \lean{IPhO2026Problems.IPhO2026_2_C_1.rayA_slope_and_intercept}
  \uses{def:physics:IPhO_2026_2_C_1:length_quantity, def:physics:IPhO_2026_2_C_1:length_in_metres, def:physics:IPhO_2026_2_C_1:aux001, def:physics:IPhO_2026_2_C_1:aux002, def:physics:IPhO_2026_2_C_1:aux003, def:physics:IPhO_2026_2_C_1:aux004, def:physics:IPhO_2026_2_C_1:aux005, def:physics:IPhO_2026_2_C_1:aux006, def:physics:IPhO_2026_2_C_1:aux007, def:physics:IPhO_2026_2_C_1:aux008, def:physics:IPhO_2026_2_C_1:aux009, def:physics:IPhO_2026_2_C_1:aux010, def:physics:IPhO_2026_2_C_1:aux011, def:physics:IPhO_2026_2_C_1:aux012, def:physics:IPhO_2026_2_C_1:aux013, def:physics:IPhO_2026_2_C_1:aux014, def:physics:IPhO_2026_2_C_1:aux015, def:physics:IPhO_2026_2_C_1:aux016, def:physics:IPhO_2026_2_C_1:aux017, lem:physics:IPhO_2026_2_C_1:aux018, lem:physics:IPhO_2026_2_C_1:aux019, lem:physics:IPhO_2026_2_C_1:aux020}
  IPhO 2026, Problem 2, C.1. In the Figure 2g coordinate convention, the reflected ray A has slope cot (2 * theta) and intercept R / (2 * cos theta).
\end{theorem}
\begin{proof}
  Reflect the upward unit direction in the radial unit normal at
  \((R\sin\theta,R\cos\theta)\).  The selected outgoing branch has components
  proportional to \((-\sin 2\theta,-\cos 2\theta)\), so the supporting-line
  slope is \(\cot(2\theta)\).  Substituting the impact point into the line
  equation and simplifying gives the SI intercept \(R/(2\cos\theta)\).
\end{proof}
\section*{Formal declaration index}
The following blocks record the typed carriers and bridge declarations used by the target.  Their dependency lists follow the declaration-level model in the covered Lean file.

\begin{definition}[Declaration LengthQuantity]
  \label{def:physics:IPhO_2026_2_C_1:length_quantity}
  \lean{IPhO2026Problems.IPhO2026_2_C_1.LengthQuantity}
  A unit-independent physical length represented by Physlib's dimensionful
  carrier with length dimension.
\end{definition}
\begin{proof}
  This is the specialization of Physlib's dimensionful-quantity interface to
  the physical dimension of length.
\end{proof}

\begin{definition}[Declaration lengthInMetres]
  \label{def:physics:IPhO_2026_2_C_1:length_in_metres}
  \lean{IPhO2026Problems.IPhO2026_2_C_1.lengthInMetres}
  \uses{def:physics:IPhO_2026_2_C_1:length_quantity}
  The scalar projection of a physical length in metres under Physlib's SI
  unit choice.
\end{definition}
\begin{proof}
  Evaluate the dimensionful length at the SI unit choice and take its
  underlying real readout.
\end{proof}

\begin{definition}[Declaration PlanePoint]
  \label{def:physics:IPhO_2026_2_C_1:aux001}
  \lean{IPhO2026Problems.IPhO2026_2_C_1.PlanePoint}
  \uses{def:physics:IPhO_2026_2_C_1:length_quantity}
  A point in the Cartesian cross-section of Figure 2g.  Both coordinates are
  Physlib length quantities; analytic line equations refer to their common
  named SI-length projection.
\end{definition}
\begin{proof}
  This is the typed definition of the stated physical carrier; its fields or defining equation provide the claimed data.
\end{proof}

\begin{definition}[Declaration PlaneDirection]
  \label{def:physics:IPhO_2026_2_C_1:aux002}
  \lean{IPhO2026Problems.IPhO2026_2_C_1.PlaneDirection}
  A dimensionless propagation direction in the Cartesian cross-section.
\end{definition}
\begin{proof}
  This is the typed definition of the stated physical carrier; its fields or defining equation provide the claimed data.
\end{proof}

\begin{definition}[Declaration dot]
  \label{def:physics:IPhO_2026_2_C_1:aux003}
  \lean{IPhO2026Problems.IPhO2026_2_C_1.PlaneDirection.dot}
  \uses{def:physics:IPhO_2026_2_C_1:aux002}
  The coordinate dot product on directions in the plane.
\end{definition}
\begin{proof}
  This is the typed definition of the stated physical carrier; its fields or defining equation provide the claimed data.
\end{proof}

\begin{definition}[Declaration IsNonzero]
  \label{def:physics:IPhO_2026_2_C_1:aux004}
  \lean{IPhO2026Problems.IPhO2026_2_C_1.PlaneDirection.IsNonzero}
  \uses{def:physics:IPhO_2026_2_C_1:aux002}
  A direction is nonzero when at least one Cartesian component is nonzero.
\end{definition}
\begin{proof}
  This is the typed definition of the stated physical carrier; its fields or defining equation provide the claimed data.
\end{proof}

\begin{definition}[Declaration IsUnit]
  \label{def:physics:IPhO_2026_2_C_1:aux005}
  \lean{IPhO2026Problems.IPhO2026_2_C_1.PlaneDirection.IsUnit}
  \uses{def:physics:IPhO_2026_2_C_1:aux002, def:physics:IPhO_2026_2_C_1:aux003}
  The normal used in the reflection formula has unit Euclidean length.
\end{definition}
\begin{proof}
  This is the typed definition of the stated physical carrier; its fields or defining equation provide the claimed data.
\end{proof}

\begin{definition}[Declaration reflectedByNormal]
  \label{def:physics:IPhO_2026_2_C_1:aux006}
  \lean{IPhO2026Problems.IPhO2026_2_C_1.PlaneDirection.reflectedByNormal}
  \uses{def:physics:IPhO_2026_2_C_1:aux002, def:physics:IPhO_2026_2_C_1:aux003}
  Reflection of a propagation direction from a surface with a unit outward normal. For an incoming direction d and normal n, this is d - 2 (d ⬝ n) n.
\end{definition}
\begin{proof}
  This is the typed definition of the stated physical carrier; its fields or defining equation provide the claimed data.
\end{proof}

\begin{definition}[Declaration IsSpecularReflection]
  \label{def:physics:IPhO_2026_2_C_1:aux007}
  \lean{IPhO2026Problems.IPhO2026_2_C_1.PlaneDirection.IsSpecularReflection}
  \uses{def:physics:IPhO_2026_2_C_1:aux002, def:physics:IPhO_2026_2_C_1:aux005, def:physics:IPhO_2026_2_C_1:aux006}
  The specular-reflection law at a mirror point, expressed as a constraining vector equation rather than as the requested slope/intercept formula.
\end{definition}
\begin{proof}
  This is the typed definition of the stated physical carrier; its fields or defining equation provide the claimed data.
\end{proof}

\begin{definition}[Declaration HalfCylindricalMirror]
  \label{def:physics:IPhO_2026_2_C_1:aux008}
  \lean{IPhO2026Problems.IPhO2026_2_C_1.HalfCylindricalMirror}
  \uses{def:physics:IPhO_2026_2_C_1:length_quantity, def:physics:IPhO_2026_2_C_1:length_in_metres}
  The upper semicircular cross-section of a half-cylindrical mirror, whose
  radius is a positive Physlib length quantity.
\end{definition}
\begin{proof}
  This is the typed definition of the stated physical carrier; its fields or defining equation provide the claimed data.
\end{proof}

\begin{definition}[Declaration OnReflectingSurface]
  \label{def:physics:IPhO_2026_2_C_1:aux009}
  \lean{IPhO2026Problems.IPhO2026_2_C_1.HalfCylindricalMirror.OnReflectingSurface}
  \uses{def:physics:IPhO_2026_2_C_1:aux001, def:physics:IPhO_2026_2_C_1:aux008}
  Membership in the reflecting upper semicircle in Figure 2g.
\end{definition}
\begin{proof}
  This is the typed definition of the stated physical carrier; its fields or defining equation provide the claimed data.
\end{proof}

\begin{definition}[Declaration SlopeInterceptLine]
  \label{def:physics:IPhO_2026_2_C_1:aux010}
  \lean{IPhO2026Problems.IPhO2026_2_C_1.SlopeInterceptLine}
  \uses{def:physics:IPhO_2026_2_C_1:length_quantity, def:physics:IPhO_2026_2_C_1:length_in_metres}
  A line in the Figure 2g convention \(y=mx+b\).  Its slope is
  dimensionless, while its intercept is a Physlib length quantity; incidence
  is stated through the coordinates' named SI projections.
\end{definition}
\begin{proof}
  This is the typed definition of the stated physical carrier; its fields or defining equation provide the claimed data.
\end{proof}

\begin{definition}[Declaration Contains]
  \label{def:physics:IPhO_2026_2_C_1:aux011}
  \lean{IPhO2026Problems.IPhO2026_2_C_1.SlopeInterceptLine.Contains}
  \uses{def:physics:IPhO_2026_2_C_1:aux001, def:physics:IPhO_2026_2_C_1:aux010}
  A point lies on a slope-intercept line.
\end{definition}
\begin{proof}
  This is the typed definition of the stated physical carrier; its fields or defining equation provide the claimed data.
\end{proof}

\begin{definition}[Declaration HasDirection]
  \label{def:physics:IPhO_2026_2_C_1:aux012}
  \lean{IPhO2026Problems.IPhO2026_2_C_1.SlopeInterceptLine.HasDirection}
  \uses{def:physics:IPhO_2026_2_C_1:aux002, def:physics:IPhO_2026_2_C_1:aux010}
  A nonvertical line has direction as a supporting direction. This equation keeps the orientation-free line geometry separate from the oriented propagation direction of an optical ray.
\end{definition}
\begin{proof}
  This is the typed definition of the stated physical carrier; its fields or defining equation provide the claimed data.
\end{proof}

\begin{definition}[Declaration OpticalRayAtPoint]
  \label{def:physics:IPhO_2026_2_C_1:aux013}
  \lean{IPhO2026Problems.IPhO2026_2_C_1.OpticalRayAtPoint}
  \uses{def:physics:IPhO_2026_2_C_1:aux001, def:physics:IPhO_2026_2_C_1:aux002, def:physics:IPhO_2026_2_C_1:aux004}
  An oriented optical ray together with a marked point on its path.
\end{definition}
\begin{proof}
  This is the typed definition of the stated physical carrier; its fields or defining equation provide the claimed data.
\end{proof}

\begin{definition}[Declaration Figure2gRayInteraction]
  \label{def:physics:IPhO_2026_2_C_1:aux014}
  \lean{IPhO2026Problems.IPhO2026_2_C_1.Figure2gRayInteraction}
  \uses{def:physics:IPhO_2026_2_C_1:aux001, def:physics:IPhO_2026_2_C_1:aux002, def:physics:IPhO_2026_2_C_1:aux007, def:physics:IPhO_2026_2_C_1:aux008, def:physics:IPhO_2026_2_C_1:aux009, def:physics:IPhO_2026_2_C_1:aux010, def:physics:IPhO_2026_2_C_1:aux011, def:physics:IPhO_2026_2_C_1:aux012, def:physics:IPhO_2026_2_C_1:aux013}
  The physical and figure-derived data for one ray striking the mirror at incidence angle incidenceAngle. Angles are measured in radians. The line coefficients remain unconstrained except through the hit-point, direction, and specular-reflection equations; in particular, neither requested closed form is assumed here.
\end{definition}
\begin{proof}
  This is the typed definition of the stated physical carrier; its fields or defining equation provide the claimed data.
\end{proof}

\begin{definition}[Declaration mA]
  \label{def:physics:IPhO_2026_2_C_1:aux015}
  \lean{IPhO2026Problems.IPhO2026_2_C_1.Figure2gRayInteraction.mA}
  \uses{def:physics:IPhO_2026_2_C_1:aux008, def:physics:IPhO_2026_2_C_1:aux014}
  The scalar readout denoted by m\_A when the interaction is ray A.
\end{definition}
\begin{proof}
  This is the typed definition of the stated physical carrier; its fields or defining equation provide the claimed data.
\end{proof}

\begin{definition}[Declaration bA]
  \label{def:physics:IPhO_2026_2_C_1:aux016}
  \lean{IPhO2026Problems.IPhO2026_2_C_1.Figure2gRayInteraction.bA}
  \uses{def:physics:IPhO_2026_2_C_1:aux008, def:physics:IPhO_2026_2_C_1:aux014}
  The Physlib length-valued intercept denoted by \(b_A\); its closed-form
  equality is stated through the named SI-length projection.
\end{definition}
\begin{proof}
  This is the typed definition of the stated physical carrier; its fields or defining equation provide the claimed data.
\end{proof}

\begin{definition}[Declaration Figure2gCausticSetup]
  \label{def:physics:IPhO_2026_2_C_1:aux017}
  \lean{IPhO2026Problems.IPhO2026_2_C_1.Figure2gCausticSetup}
  \uses{def:physics:IPhO_2026_2_C_1:aux001, def:physics:IPhO_2026_2_C_1:aux008, def:physics:IPhO_2026_2_C_1:aux011, def:physics:IPhO_2026_2_C_1:aux014}
  Context retained for the later caustic construction: ray B is parallel to ray A, has incidence angle theta + deltaTheta, and intersects ray A's reflected line at the sampled caustic point. relativeScale makes the informal condition Δθ ≪ θ explicit without choosing an unstated numerical tolerance.
\end{definition}
\begin{proof}
  This is the typed definition of the stated physical carrier; its fields or defining equation provide the claimed data.
\end{proof}

\begin{lemma}[Declaration reflected\_direction\_from\_specular\_law]
  \label{lem:physics:IPhO_2026_2_C_1:aux018}
  \lean{IPhO2026Problems.IPhO2026_2_C_1.reflected_direction_from_specular_law}
  \uses{def:physics:IPhO_2026_2_C_1:aux008, def:physics:IPhO_2026_2_C_1:aux014}
  The vector reflection law and Figure 2g orientation select the outgoing down-left branch and give its doubled-angle direction.
\end{lemma}
\begin{proof}
  Substitute the cited governing relations in order and use the stated positivity, orientation, and branch conditions to obtain the conclusion.
\end{proof}

\begin{lemma}[Declaration reflected\_line\_slope]
  \label{lem:physics:IPhO_2026_2_C_1:aux019}
  \lean{IPhO2026Problems.IPhO2026_2_C_1.reflected_line_slope}
  \uses{def:physics:IPhO_2026_2_C_1:aux008, def:physics:IPhO_2026_2_C_1:aux014, def:physics:IPhO_2026_2_C_1:aux015, lem:physics:IPhO_2026_2_C_1:aux018}
  The direction equation of the reflected supporting line determines its dimensionless slope.
\end{lemma}
\begin{proof}
  Substitute the cited governing relations in order and use the stated positivity, orientation, and branch conditions to obtain the conclusion.
\end{proof}

\begin{lemma}[Declaration reflected\_line\_intercept]
  \label{lem:physics:IPhO_2026_2_C_1:aux020}
  \lean{IPhO2026Problems.IPhO2026_2_C_1.reflected_line_intercept}
  \uses{def:physics:IPhO_2026_2_C_1:aux008, def:physics:IPhO_2026_2_C_1:aux014, def:physics:IPhO_2026_2_C_1:aux016, lem:physics:IPhO_2026_2_C_1:aux018, lem:physics:IPhO_2026_2_C_1:aux019}
  Incidence of the reflected line at the mirror hit point determines its length-valued intercept.
\end{lemma}
\begin{proof}
  Substitute the cited governing relations in order and use the stated positivity, orientation, and branch conditions to obtain the conclusion.
\end{proof}
% --- Archon physics formalization source end ---
